\section{Two ways to describe a type}
\label{sec:ch02-two-ways}
\index{schema!authoring styles|(}

Every GraphQL server in .NET has to answer the same question: where does the
schema come from? You can write C\# and let the server work out the schema, or
you can write the schema and let the server work out where to send the calls.
HotChocolate has historically supported both, in three named flavours, and if
you have read about this before you probably learned them as annotation-based,
code-first, and schema-first.

That vocabulary is out of date. HotChocolate~16's documentation names two
authoring styles under a heading that says so plainly: \enquote{Hot Chocolate
offers two C\# authoring styles, both of which produce GraphQL
SDL}~\autocite{chillicream2026typesystem}. They are called
\emph{implementation-first} and \emph{code-first}. The term
\enquote{annotation-based} has been retired, and the third style has quietly
left the documentation altogether. It still works; it is just no longer written
down.

The comparison below uses one deliberately small schema so the three versions
fit side by side. It lives in the companion repository under
\code{samples/three-approaches}, as three projects that run on three ports. All
three serve exactly this:

\begin{graphqlsdl}
type Query {
  productById(id: ID!): Product
}

type Product {
  id: ID!
  sku: String!
  title: String!
}
\end{graphqlsdl}

\subsection*{Implementation-first}

\index{schema!implementation-first}%
C\# types are the source of truth, and attributes adjust what the server infers
from them.

\begin{minted}{csharp}
[QueryType]
public static partial class CatalogQueries
{
    public static Product? GetProductById([ID] Guid id) => Catalog.ById(id);
}

[ObjectType<Product>]
public static partial class ProductNode
{
    [ID]
    public static Guid GetId([Parent] Product product) => product.Id;
}
\end{minted}

\code{Product} itself is a plain record, and its \code{Sku} and \code{Title}
properties become fields with no further ceremony. The only reason
\code{ProductNode} exists is to say that the identifier is an \code{ID} and not
a \code{UUID}, which is the default mapping for a \code{Guid}. Registration is
the one generated line from the previous section:

\begin{minted}{csharp}
builder.AddGraphQL().AddCatalog();
\end{minted}

\subsection*{Code-first}

\index{schema!code-first}%
The same schema, described through the descriptor API. Nothing is inferred.

\begin{minted}{csharp}
public sealed class ProductType : ObjectType<Product>
{
    protected override void Configure(IObjectTypeDescriptor<Product> descriptor)
    {
        descriptor.Name("Product");

        descriptor.Field(p => p.Id).Type<NonNullType<IdType>>();
        descriptor.Field(p => p.Sku).Type<NonNullType<StringType>>();
        descriptor.Field(p => p.Title).Type<NonNullType<StringType>>();
    }
}
\end{minted}

Registration names every type, because a type that is not reachable from
\code{AddQueryType} through \code{AddType} is not in the schema:

\begin{minted}{csharp}
builder.AddGraphQL()
    .AddQueryType<QueryType>()
    .AddType<ProductType>();
\end{minted}

It is worth saying that this API survived HotChocolate~16 untouched. The release
rearchitected the type system, and I half expected the descriptors to change
shape along with it, but \code{ObjectType<T>},
\code{IObjectTypeDescriptor<T>} and \code{Configure} are exactly where they were.
Code-first is not a legacy path.

\subsection*{And schema-first}

\index{schema!schema-first}%
The third style writes SDL and binds runtime types to it.

\begin{minted}{csharp}
const string Sdl =
    """
    type Query {
      productById(id: ID!): Product
    }

    type Product {
      id: ID!
      sku: String!
      title: String!
    }
    """;

builder.AddGraphQL()
    .AddDocumentFromString(Sdl)
    .BindRuntimeType<Product>("Product")
    .AddResolver<QueryResolvers>("Query");
\end{minted}

That compiles and runs against 16.6.0, and I have the exported schema to prove
it. But there is no page about it anywhere in the version~16 documentation. Not
a deprecated page, not a page with a warning banner: the whole documentation
tree has nothing on schema-first, and the tab that used to sit alongside the
other two in the code samples is gone.

The API is alive, though. All four of the calls above ship in 16.6.0 without an
\code{[Obsolete]} attribute between them, and the repository still has tests
exercising them at that tag~\autocite{chillicream2026source}. So the honest
description is: supported, tested, and undocumented. Whether that is a
deliberate soft retirement or an oversight I cannot tell you, and I would
rather say so than guess.

Two things have moved, if you are porting older schema-first code.
\code{BindComplexType} is gone and \code{BindRuntimeType} replaces it, and
\code{AddSchemaConfiguration} is now internal and cannot be called from your
code. Search results will happily hand you pages describing both, from a
documentation repository ChilliCream retired years ago. Stale documentation
outranking current documentation is not specific to HotChocolate, and the
defence is the same everywhere: for anything version-sensitive, read the source
tree at a release tag before you read the first page of results.

\subsection*{What Mosaic uses, and why}

Mosaic is implementation-first throughout, and since every chapter from here
inherits that decision it deserves a reason rather than a preference.

\index{HotChocolate!analyzers}%
The strongest argument is that HotChocolate~16 has put its tooling behind it.
The source generator, the \code{[DataLoader]} attribute that
chapter~\ref{ch:data-without-the-n-plus-1} needs, and a set of compile-time
analyzers with code fixes all key off the attributes; version~16 added twelve
new analyzers at error severity on its own~\autocite{chillicream2026migrate16}.
Descriptors get none of that: a mistake in a \code{Configure} method is a
runtime schema error, while the same mistake in an attributed class is a red
squiggle with a suggested fix. For a book whose thesis is that you should earn
the right to distribute before you distribute, taking the compile-time help is
the consistent choice.

Second, the generated code is something we can open and read.
Chapter~\ref{ch:the-life-of-a-request} answers the question \enquote{how does
HotChocolate know about my types} by looking at the file the generator wrote,
and chapter~\ref{ch:inside-hotchocolate} goes considerably further. Descriptors
would leave nothing to look at but reflection.

Third, and this is the one that matters most for the shape of the book, an
extension class lets a field live in a different file from the type it hangs
off, with no central registration list to edit. That is what
section~\ref{sec:ch02-six-domains} is about, and it is what makes
chapter~\ref{ch:the-first-cut-extracting-catalog} a move rather than a rewrite.

Fourth, and last, federation is expressed with attributes. Part~II runs on
\code{[Key]} and \code{[ReferenceResolver]}, which are the documented path in
the federation package, so an implementation-first service is already speaking
the right language when it becomes a subgraph.

I want to be precise about what that decision is not. It is not a claim that
descriptors are worse. They are the escape hatch, and they are the only route to
a schema whose shape is decided at runtime. Mosaic will reach for one when an
attribute cannot express what a type needs, and when that happens in Part~III it
should read as planned rather than as a reversal.

\subsection*{One schema, three descriptions}

Run all three sample projects and export all three schemas, and the files are
identical. Not equivalent: identical, same bytes, same hash. The style you pick
changes who maintains what and which mistakes the compiler catches. It changes
nothing at all about what a client sees.

\begin{figure}[htbp]
  \centering
  % Three ways to describe a type, one schema, one SDL document. Referenced from
% chapter 02. Bare tikzpicture: the figure environment, caption and label live
% at the call site. Uses only arrows.meta and positioning (loaded in packages.tex).
\begin{tikzpicture}[
    font=\small,
    src/.style={draw, rounded corners=2pt, minimum width=40mm,
                minimum height=17mm, align=center, font=\scriptsize},
    odd/.style={draw, dashed, rounded corners=2pt, minimum width=40mm,
                minimum height=17mm, align=center, font=\scriptsize},
    core/.style={draw, thick, rounded corners=2pt, minimum width=34mm,
                 minimum height=12mm, align=center},
    emit/.style={draw, rounded corners=2pt, minimum width=32mm,
                 minimum height=12mm, align=center},
    flow/.style={-{Stealth[length=2mm]}, thick},
    note/.style={font=\scriptsize\itshape, align=center}
  ]

  \node[src] (impl) at (-4.8,2.5)
    {{\small Implementation-first}\\[2pt]
     C\# types plus attributes,\\
     read by a source generator};

  \node[src] (code) at (-4.8,0)
    {{\small Code-first}\\[2pt]
     descriptor classes:\\
     \texttt{ObjectType<T>} and\\
     \texttt{Configure}};

  \node[odd] (schemafirst) at (-4.8,-2.5)
    {{\small Schema-first}\\[2pt]
     an SDL document plus\\
     runtime type bindings};

  \node[core] (schema) at (0,0)   {one schema object\\in memory};
  \node[emit] (sdl)    at (4.5,0) {the same SDL,\\byte for byte};

  \foreach \s in {impl,code,schemafirst}{%
    \draw[flow] (\s) -- (schema);
  }
  \draw[flow] (schema) -- (sdl);

  \node[note, anchor=south] at (-4.8,3.45)
    {Mosaic uses implementation-first};
  \node[note, anchor=north] at (-4.8,-3.5)
    {dashed: supported, but absent\\from the v16 documentation};

\end{tikzpicture}

  \caption{Three descriptions, one schema, one SDL document. The choice is a
  maintenance decision made inside the service, and it is invisible from
  outside it. Schema-first is drawn apart because it still works but no longer
  appears in the version~16 documentation.}
  \label{fig:ch02-authoring-styles}
\end{figure}

That is also why the exercise at the end of this chapter asks you to prove it
rather than take my word for it.
\index{schema!authoring styles|)}
